\documentclass[11pt,letterpaper]{article}
\usepackage[margin=.83in,top=.72in,bottom=.75in]{geometry}
\usepackage{amsmath,amssymb,booktabs,array,xcolor,fancyhdr,microtype,hyperref}
\definecolor{ink}{HTML}{122843}\definecolor{accent}{HTML}{176865}
\hypersetup{colorlinks=true,urlcolor=accent,linkcolor=accent,pdftitle={Exact periods of two central-binomial Stirling transforms}}
\pagestyle{fancy}\fancyhf{}\renewcommand{\headrulewidth}{0pt}
\fancyfoot[L]{\footnotesize\color{accent}Exact modular periods of A354242 and A354253}\fancyfoot[R]{\footnotesize\color{accent}\thepage}
\setlength{\parindent}{0pt}\setlength{\parskip}{5pt}\setlength{\footskip}{24pt}
\newcommand{\head}[1]{\vspace{6pt}{\large\sffamily\bfseries\color{ink}#1}\par}
\newcommand{\overlinehead}[1]{{\footnotesize\sffamily\bfseries\color{accent}#1}\par}
\newcommand{\Fp}{\mathbb F_p}\newcommand{\Z}{\mathbb Z}\newcommand{\lcm}{\operatorname{lcm}}
\newcommand{\oeis}[1]{\href{https://oeis.org/#1}{#1}}
\begin{document}
\overlinehead{NUMBER THEORY \quad / \quad RESEARCH NOTE \hfill 26 NOVEMBER 2026}
\vspace{4pt}{\fontsize{24.88}{28}\selectfont\sffamily\bfseries\color{ink}Exact periods of two central-binomial\\ Stirling transforms}\par
For $b=1,2$ define the integer sequence $a_b(n)$ by its exponential generating function
\[
 \sum_{n\ge0} a_b(n)\frac{x^n}{n!}=(1-4b(e^x-1))^{-1/2}.
\]
These are \oeis{A354242} and \oeis{A354253}, respectively. Peter Bala conjectured that the least eventual period modulo $m$ divides Euler's totient $\varphi(m)$. An earlier proof treats the whole class $G(e^x-1)$ with integral $G$ and gives $\varphi(m)$ as a period from $n\ge m$ [3]. We determine the least eventual period of each of these two particular sequences for every positive modulus.

\head{The exact answer}
Let $T_b(m)$ be the least positive $T$ such that $a_b(n+T)\equiv a_b(n)\pmod m$ for every sufficiently large $n$. For $e\ge1$ the prime-power answers are
\begin{center}\renewcommand{\arraystretch}{1.33}
\begin{tabular}{@{}lll@{}}\toprule
\textbf{Modulus} & \textbf{$T_1$, A354242}& \textbf{$T_2$, A354253}\\\midrule
$p^e$, prime $p\ge5$ & $(p-1)p^{e-1}$ & $(p-1)p^{e-1}$\\
$3^e$ & $1$ if $e=1$ & $3^{e-1}$ if $e\le5$\\
 & $2\cdot3^{e-1}$ if $e\ge2$ & $2\cdot3^{e-1}$ if $e\ge6$\\
$2^e$ & $1$ if $e\le3$ & $1$ if $e\le6$\\
 & $2^{\max(1,e-5)}$ if $e\ge4$ & $2^{\max(1,e-7)}$ if $e\ge7$\\\bottomrule
\end{tabular}\end{center}
For $m>1$, take the least common multiple over its exact prime-power factors. Set $T_b(1)=1$. This follows from the Chinese remainder theorem, since a shift works eventually modulo $m$ exactly when it works eventually at every factor. The eventual shifts at a fixed modulus are precisely the multiples of its least eventual period, by division with remainder on a repeating tail.

\head{Two examples}
Modulo $243=3^5$, the two periods are $162$ and $81$. At the very next power, $729$, both periods are $486$. Modulo $1024=2^{10}$ they are $32$ and $8$. The distinction already appears modulo $64$, where the first sequence eventually alternates between two residues and the second becomes constant.

\newpage
\overlinehead{PROOF \quad / \quad REDUCTION TO FINITE SUMS}
\head{A finite sum of integer powers at every modulus}
Write $g_k=b^k\binom{2k}{k}$. Expanding the generating function in powers of $e^x-1$ gives
\[
 a_b(n)=\sum_{k=0}^n g_k\,k!\,\genfrac{\{}{\}}{0pt}{}{n}{k},\qquad
 k!\,\genfrac{\{}{\}}{0pt}{}{n}{k}=\sum_{j=0}^k(-1)^{k-j}\binom{k}{j}j^n.
\]
Here $\genfrac{\{}{\}}{0pt}{}{n}{k}$ is a Stirling number of the second kind, and $0^0=1$. The second identity is inclusion--exclusion for onto maps and holds also for $k>n$, when both sides vanish. For a modulus $m>1$, take any $K$ with $m\mid K!$. Define integer weights by
\[
 P_K(y)=\sum_{k=0}^{K-1}g_k(y-1)^k=\sum_{j=0}^{K-1}c_jy^j.
 \qquad a_b(n)\equiv\sum_{j=0}^{K-1}c_jj^n\pmod m. \tag{1}
\]
Every omitted Stirling term is a multiple of $m$. Modulo $p^e$, all bases divisible by $p$ vanish once $n\ge e$. The remaining bases are units, so the Carmichael exponent $\lambda(p^e)$ is a period on that tail. A value contradicting a proposed shift anywhere on this tail recurs arbitrarily late.

\head{The nonzero bases that control the period modulo $p$}
For odd $p$ write $h=(p-1)/2$ and $H(y)=(1+4b-4by)^h=\sum_{r=0}^h w_r y^r$ in $\Fp[y]$. Choose $K\ge p^2$ in (1), increasing it further if the modulus requires it. In $\Fp[y]$ one has
\[
 P_K(y)\equiv H(y)\pmod{y^p-1},\qquad
 P_K(y)\equiv H(y)H(y^p)\pmod{y^{p^2}-1}. \tag{2}
\]
Indeed in $\Fp[[t]]/(t^q)$, for $q=p,p^2$, $(1-4bt)^{(q-1)/2}$ and $(1-4bt)^{-1/2}$ have the same square and constant coefficient $1$, since $(1-4bt)^q=1$ there. Their sum is a unit, so they agree. Substitute $t=y-1$ and use $(y-1)^q=y^q-1$. For $q=p^2$ the exponent is $h(1+p)$, which gives the factorization in (2).

For $n\ge1$ the reduction modulo $p$ is $\sum_{r=1}^h w_r r^n$. Distinct nonzero bases give linearly independent sequences on every tail. Evaluating at consecutive exponents produces a Vandermonde matrix multiplied by an invertible diagonal matrix. A shift $d$ therefore works exactly when $r^d=1$ for every $r$ of nonzero weight. If $p\nmid b(1+4b)$ all the weights are nonzero. For $p\ge5$ the bases $1,\ldots,h$ generate $\Fp^\times$, since $2h=-1$ modulo $p$ and their negatives cover the other half. The least period is $p-1$. If $p\mid1+4b$, the polynomial is $y^h$. This gives period $4$ for $(b,p)=(1,5)$ and period $1$ for $(2,3)$. For $(1,3)$ the only nonzero base is also $1$.

\head{Sharp universal bounds for the whole class}
For any integral series $G$, put $a_G(n)=n![x^n]G(e^x-1)$. For $m>1$ let $E(m)$ be the largest exponent in the prime factorization of $m$, and let $\lambda(m)$ be the exponent of its unit group. For $n\ge0$ and $d>0$ we have the sharp equivalence
\[
 a_G(n+d)\equiv a_G(n)\pmod m\ \text{for every }G
 \quad\Longleftrightarrow\quad n\ge E(m)\ \text{and}\ \lambda(m)\mid d.
\]
Formula (1) works with the coefficients of any $G$. For sufficiency, reduce at each $p^e\Vert m$. A nonunit base has zero $n$th power there, and every unit has $d$th power $1$. For necessity, $G(t)=(1+t)^j$ gives $a_G(n)=j^n$. Testing every unit $j$ forces $\lambda(m)\mid d$. Taking $j=p$ at a factor $p^{E(m)}\Vert m$ forces $n\ge E(m)$, since $p\nmid p^d-1$.

\newpage
\overlinehead{PROOF \quad / \quad ODD PRIME POWERS}
\head{Why the full power of an odd prime occurs}
Fix odd $p$ and $e\ge2$ and put $D=(p-1)p^{e-2}$. For $p\nmid j$ set $q_p(j)=(j^{p-1}-1)/p$. The binomial theorem yields
\[
 j^D\equiv 1+p^{e-1}q_p(j)\pmod {p^e}.
\]
Starting at $e=2$, raise the congruence to the $p$th power at each step. The intermediate binomial coefficients are multiples of $p$ and the last term vanishes to the required order because $p$ is odd. Also, for $1\le r<p$, expansion modulo $p^2$ gives
\[
 q_p(r+p\ell)\equiv q_p(r)-\ell/r\pmod p.
\]
For $n\ge e$ use (1) modulo $p^e$, divide $a_b(n+D)-a_b(n)$ by $p^{e-1}$ and reduce modulo $p$. Only unit bases contribute. In (2), the coefficient at $y^{r+p\ell}$ is the sum of the $c_j$ in that residue class modulo $p^2$, reduced modulo $p$. This coefficient is $w_rw_\ell$ (take $w_i=0$ for $i>h$). Because $\sum_\ell w_\ell=H(1)=1$ and $\sum_\ell\ell w_\ell=H'(1)=-4bh=2b$ in $\Fp$, we get the exact test
\[
 \frac{a_b(n+D)-a_b(n)}{p^{e-1}}
 \equiv\sum_{r=1}^{h} w_r\big(q_p(r)-2b/r\big)r^n\pmod p. \tag{3}
\]
If $p\nmid b(1+4b)$ the coefficient of $1^n$ is $-2b w_1\ne0$. For $(b,p)=(1,5)$ the surviving base is $r=2$ and $q_5(2)-2b/r=3-1=2$ modulo $5$. For $(2,3)$ the surviving base is $r=1$ and $q_3(1)-2b=-4\ne0$ modulo $3$. Thus (3) is nonzero somewhere on every tail by the same Vandermonde argument. The least period divides $(p-1)p^{e-1}$ and cannot divide $D$. Its $p$-part is exactly $p^{e-1}$. For $p\ge5$ its reduction modulo $p$ also forces the factor $p-1$. This proves that entire row of the table.

\head{When the factor 2 appears at the prime 3}
For $b=1$, the residues $a_1(2)=5$ and $a_1(5)=2$ modulo 9 exclude shift $3$ on the periodic tail. The least period modulo 9 is $6$, so every higher period is even. For $b=2$, shift $81$ works modulo $243$ from $n=5$. It follows that the periods at $3^e$ for $e\le5$ are odd. Modulo $729$, $a_2(6)=226$ and $a_2(249)=712$ exclude shift $243$, so the period is even from $e=6$ on. Together with (3) this proves both ternary rows. Page 5 supplies the entire finite certificate for the positive assertion at $243$.

\newpage
\overlinehead{PROOF \quad / \quad POWERS OF TWO}
\head{One doubling calculation controls arbitrarily large exponents}
Put $\Delta_s(n)=a_b(n+2^s)-a_b(n)$. Let $v_b(s)$ be the largest $v$ for which $2^v\mid\Delta_s(n)$ eventually. We will show
\[
\begin{array}{c|cc}
 &s=0&s\ge1\\ \hline
 b=1 &3&s+5\\
 b=2 &6&s+7
\end{array} \tag{4}
\]
For $s\ge1$ write $c=5$ for $b=1$ and $c=7$ for $b=2$, and put $R=c-2$. The finite checks on page 5 establish (4) through $s=R$. For an odd integer $j$, define $L_s(j)=(j^{2^s}-1)/2^{s+2}$. These are integers for $s\ge1$, since $j^2\equiv1\pmod8$ and squaring gives
\[
 L_{s+1}(j)=L_s(j)+2^{s+1}L_s(j)^2.
\]
Consequently $L_s(j)\equiv L_R(j)\pmod{2^{c-1}}$ whenever $s\ge R$. Multiplying by $2^{s+2}$ gives
\[
 j^{2^s}-1\equiv 2^{s-R}(j^{2^R}-1)\pmod {2^{s+c+1}}. \tag{5}
\]
Use the \emph{same} weights from (1) modulo $2^{s+c+1}$ on both sides. For even $j$, $j^n$ already vanishes when $n\ge s+c+1$. For odd $j$, multiply (5) by $c_j j^n$ and sum. Hence
\[
 \Delta_s(n)\equiv 2^{s-R}\Delta_R(n)\pmod {2^{s+c+1}}\qquad(n\ge s+c+1).
\]
The anchor $\Delta_R$ is eventually divisible by $2^{R+c}$ and attains exactly this valuation arbitrarily late. The congruence gives the same two claims at valuation $s+c$ for $\Delta_s$. This proves (4) for every $s$.

The least eventual period modulo $2^e$ is a power of two, because it divides $\lambda(2^e)$. By (4) it is precisely $2^s$ for the smallest nonnegative $s$ with $v_b(s)\ge e$. This yields the two dyadic formulas on page 1.

\head{How a finite window proves an infinite tail}
For any shift $d$, (1) writes the difference as a sum of $K$ terms $c_j(j^d-1)j^n$. It satisfies the monic linear recurrence with characteristic polynomial $\prod_{j=0}^{K-1}(X-j)$ over $\Z/m\Z$. Thus $K$ consecutive zeros imply zero at every later index. This use of a monic recurrence also works after reducing to a divisor of $m$. At a prime-power modulus, each nonzero window value at an index $n\ge e$ repeats arbitrarily late by the unit-base argument on page 2.

\newpage
\overlinehead{CERTIFICATES \quad / \quad REPRODUCIBILITY}
\head{Thirteen small windows}
For the dyadic checks take $m=2^{v+1}$, $N=v+1$ and $d=2^s$. In each row every difference from $n=N$ through $N+K-1$ is $0$ or $2^v$ modulo $m$. The last column supplies one index where it is $2^v$. Formula (1) with the given $K$ verifies the full window using integer arithmetic.
\begin{center}\renewcommand{\arraystretch}{1.05}
\begin{tabular}{@{}rrrrr@{\hspace{.48in}}rrrrr@{}}\toprule
$b$&$s$&$v$&$K$&Witness $n$&$b$&$s$&$v$&$K$&Witness $n$\\\midrule
1&0&3&6&4&2&0&6&8&7\\ 1&1&6&8&7&2&1&8&12&9\\
1&2&7&10&9&2&2&9&12&10\\ 1&3&8&12&9&2&3&10&14&11\\
&&&&&2&4&11&16&12\\ &&&&&2&5&12&16&13\\\bottomrule\end{tabular}\end{center}
For the three ternary checks, the parameters are $(b,m,N,d,K)$. The differences are reduced to $[0,m)$ and written in increasing index order.
\begin{center}\renewcommand{\arraystretch}{1.13}
\begin{tabular}{@{}ll@{}}\toprule Parameters & Complete difference window\\\midrule
$(1,9,2,3,6)$ & $(6,3)$ repeated three times\\
$(2,243,5,81,12)$ & twelve zeros\\
$(2,729,6,243,15)$ & $(486,243)$ repeated seven times, followed by $486$\\\bottomrule\end{tabular}\end{center}
\head{Independent checks}
The Python verifier also generates the original coefficients from $(1+4b-4be^x)F'_b(x)=2be^xF_b(x)$. This recurrence agrees with the Stirling calculation through index $100$ for each sequence. Two full Carmichael periods were compared for every modulus $1$ to $512$, and also for $729$, $1000$ and $1024$, giving 1,030 independently measured periods. The proof covers every positive modulus.

Lean checks all 13 finite windows, their propagation over commutative rings, the Vandermonde period criterion, and the integer lifting identities at odd primes and at $2$. The generating-function reductions (1)--(2), the derivation of (3) from the weights, and the assembly of the classification are proved in this note. The Python verifier also checks (3) in 888 cases. The \href{https://github.com/DomTheDeveloper/open-math-agents/tree/research/exact-stirling-periods/research/exact-stirling-periods}{source files and verification records} accompany the proof.

{\small \textbf{Sources.} [1] OEIS Foundation, \oeis{A354242}, including Peter Bala's conjecture of July 7, 2022. [2] OEIS Foundation, \oeis{A354253}, July 8, 2022. [3] The Omega Institute, \href{https://the-omega-institute.github.io/trureturing-mdbook/Blueprint/D5/S1/Recurrence/Periodic/IntegralEgfTotientPeriod.html}{\emph{Totient Periods for Integral Exponential Substitutions}}, prior proof of the general period bound. Historical priority for the minimum-period formulas is unverified.}
\end{document}
